\documentclass[11pt, a4paper, oneside]{scrartcl}

% --- Pakete für Sprache und Zeichenkodierung ---
\usepackage[utf8]{inputenc}
\usepackage[T1]{fontenc}
\usepackage[ngerman]{babel}
\usepackage[german=swiss]{csquotes}

% --- Layout und Design ---
\usepackage[margin=2.5cm]{geometry}
\usepackage{charter} % Eine gut lesbare Serifenschrift für Dokumente
\usepackage{microtype}
\usepackage{enumitem}
\usepackage{booktabs}

% --- Hyperlinks ---
\usepackage[colorlinks=true, linkcolor=blue, urlcolor=cyan]{hyperref}

% --- Dokumentmetadaten ---
\title{Expert Briefing Document}
\subtitle{Strukturvalidierung eines systemdynamischen Modells zur Landnutzungsdynamik in den EU27 (2026--2050)}
\author{Modellierungsteam}
\date{\today}

\begin{document}

\maketitle
\newpage

\section{Zielsetzung der Anfrage}
Im Rahmen eines systemdynamischen Modellierungsprojekts untersuchen wir die strukturelle Wechselwirkung zwischen Urbanisierungsprozessen und landwirtschaftlicher Flächennutzung in den EU27 im Zeitraum 2026--2050.

Ziel dieser Anfrage ist \textbf{nicht} die Validierung von Simulationsergebnissen oder Parametrisierungen, sondern die \textbf{Strukturvalidierung der angenommenen kausalen Mechanismen} auf konzeptioneller Ebene (Causal Loop Diagram). 

Wir möchten prüfen,
\begin{itemize}
  \item Fachliche Plausibilität der modellierten Einflussbeziehungen.
  \item Identifikation fehlender relevanter Mechanismen.
  \item Konsistenz der Wirkungsrichtungen mit dem Stand der Fachliteratur.
  \item Ob zentrale Rückkopplungsstrukturen adäquat abgebildet sind. 
\end{itemize}

\section{Modellrahmen und Systemgrenzen}
Das Modell aggregiert die EU27 als makroskopisches Gesamtsystem. Regionale Differenzierungen werden in dieser Phase bewusst nicht modelliert.

\paragraph{Simulationszeitraum:} 2026--2050. Historische Daten (2000--2025) dienen der konzeptionellen Kalibrierung.

\paragraph{Ausschlüsse (Out of Scope):}
\begin{itemize}
  \item Explizite Klimadynamiken und geopolitische Schocks.
  \item Globale Marktgleichgewichte.
  \item Detaillierte sektorale Differenzierungen.
\end{itemize}

Die Modellierung folgt einem strukturellen Ansatz im Sinne der \textit{System Dynamics} (Forrester), mit Fokus auf endogene Rückkopplungen.

\section{Konzeptionelle Kernmechanismen}
Basierend auf dem Causal Loop Diagram (CLD) wurden vier zentrale strukturelle Dynamiken identifiziert:

\subsection{Wirtschaftlicher Verstärkungsmechanismus}
Wirtschaftliches Wachstum erhöht das verfügbare Einkommen, was den Flächenverbrauch pro Kopf (Wohnraum, Gewerbe) steigert. Dies führt zu einer erhöhten Landumwandlung (\textit{Land Conversion}) und reduziert die landwirtschaftliche Nutzfläche.
Ohne regulatorische Gegenmechanismen wirkt dieser Prozess selbstverstärkend.

\subsection{Technologischer Kompensationsmechanismus}
Agrartechnologischer Fortschritt steigert den Ertrag pro Hektar. Bei gleichbleibender Produktionsnachfrage sinkt dadurch der theoretische Flächenbedarf. 
Dieser Mechanismus wirkt stabilisierend (Balancing Loop) auf die Flächenverfügbarkeit, kann jedoch Landverluste durch Urbanisierung oft nur partiell kompensieren.

\subsection{Raumplanerischer Schutzmechanismus}
Institutionelle Eingriffe (Spatial Planning, Green Belts) stärken den Bodenschutz. Dieser Mechanismus ist politisch sensitiv und agiert als dämpfender Faktor auf den Urbanisierungsdruck, 
ist jedoch nicht automatisch systemisch dominant.

\subsection{Importdruck-Mechanismusi}
Eine sinkende Nahrungsmittelproduktion erhöht die Importabhängigkeit. 
Der daraus resultierende wirtschaftliche und politische Druck kann wiederum Flächennutzungsentscheidungen (z.\,B. Re-Aktivierung von Brachen) beeinflussen.
Dieser Mechanismus stellt eine indirekte Rückkopplung dar.

\section{Zentrale Annahmen}
\begin{itemize}
  \item \textbf{Aggregation:} Die Aggregation über die EU27 wird als strukturell zulässig erachtet.
  \item \textbf{Pfadabhängigkeit:} \textit{Land Conversion} beschreibt primär die Umwandlung von Agrarfläche in Siedlungsfläche.
  \item \textbf{Exogenität:} Die Bevölkerungsentwicklung wirkt als exogene Treibkraft auf die Wohnraumnachfrage.
  \item \textbf{Monotonie:} Beziehungen werden als monoton gerichtet modelliert.
\end{itemize}

Diese Annahmen sind explizit dokumentiert und offen zur Diskussion.

\section{Konkrete Fragen an die Experten:innen}
Um die Validierung zielgerichtet zu gestalten, bitten wir um Stellungnahme zu folgenden Punkten:

\begin{enumerate}
  \item Sind die identifizierten Kernmechanismen relevant und korrekt? 
  \item Sind die identifizierten Kernmechanismen vollständig oder fehlen strukturell relevante Einflussfaktoren?
  \item Sind die angenommenen Wirkungsrichtungen fachlich konsistent?
  \item Gibt es bekannte Nichtlinearitäten oder Schwellenwerte, die strukturell berücksichtigt werden sollten?
  \item Sind relevante zeitliche Verzögerungen konzeptionell falsch eingeschätzt?
  \item Gibt es Literatur, die einzelne Mechanismen explizit stützt oder infrage stellt? (eher streichen - unser Job I guess)
\end{enumerate}

\section{Beilage}
\begin{itemize}
  \item Causal Loop Diagram (CLD) Grafik.
  \item Strukturierte Dokumentation.
  \item Detaillierte Loop-Beschreibungen.
\end{itemize}

\end{document}
